\section{The \est{} Estimator}
\label{sec:estimator}

\subsection{Latent-variable formulation}

Stack the received symbols as $\vect{y} \in \mathbb{C}^{N}$ and collect the
phase samples as $\boldsymbol{\phi} \in \mathbb{R}^{N}$. Using
Eq.~\eqref{eq:time-domain} the likelihood factorises as
\begin{equation}
  \label{eq:likelihood}
  p(\vect{y} \mid \vect{h}, \boldsymbol{\phi})
  = \prod_{n=0}^{N-1}
    \mathcal{CN}\!\left( y_n ;\;
      e^{j\phi_n} \bigl[\mat{X} \vect{h}\bigr]_n , \; \sigma^2 \right) ,
\end{equation}
where $\mat{X}$ is the circulant convolution matrix of the pilot block. The
phase prior is the Wiener process of Section~\ref{sec:system-model},
\begin{equation}
  \label{eq:phase-prior}
  p(\boldsymbol{\phi}) = p(\phi_0)
    \prod_{n=1}^{N-1} \mathcal{N}\!\left( \phi_n ;\, \phi_{n-1},\,
      2\pi \beta T_s \right) ,
\end{equation}
and the channel prior is $\vect{h} \sim \mathcal{CN}(\vect{0}, \mat{R}_h)$ with
$\mat{R}_h$ the exponential power delay profile of the deployment. Exact
inference in this model is intractable because $\phi_n$ enters through a
complex exponential.

\subsection{Variational updates}

Variational treatments of joint synchronisation and detection have precedent~\citep{whitfield2017variational}, though not with an exact leakage covariance. We adopt a factorised posterior
$q(\vect{h}, \boldsymbol{\phi}) = q(\vect{h}) \, q(\boldsymbol{\phi})$ and
minimise the Kullback-Leibler divergence to the true posterior. The resulting
coordinate updates are the content of the method, and both have closed forms.

Holding $q(\boldsymbol{\phi})$ fixed, the channel update is a linear minimum
mean square error solution against a phase-derotated observation. Writing
$\alpha_n = \mathbb{E}_q[e^{-j\phi_n}]$ and
$\mat{A} = \diag(\alpha_0, \dots, \alpha_{N-1})$, we have
\begin{align}
  \label{eq:h-update-1}
  q^{\star}(\vect{h})
    &\propto \exp\Bigl(
      \mathbb{E}_{q(\boldsymbol{\phi})}\bigl[ \log p(\vect{y} \mid
        \vect{h}, \boldsymbol{\phi}) \bigr]
      + \log p(\vect{h}) \Bigr) \\
  \label{eq:h-update-2}
    &= \mathcal{CN}\!\left( \vect{h} ; \, \hat{\vect{h}}, \, \boldsymbol{\Sigma}_h
       \right) , \\
  \label{eq:h-update-3}
  \boldsymbol{\Sigma}_h
    &= \left( \sigma^{-2} \mat{X}^{\mathsf{H}} \mat{D} \mat{X}
       + \mat{R}_h^{-1} \right)^{-1} , \\
  \label{eq:h-update-4}
  \hat{\vect{h}}
    &= \sigma^{-2} \boldsymbol{\Sigma}_h \mat{X}^{\mathsf{H}} \mat{A} \vect{y} ,
\end{align}
where $\mat{D} = \diag(\mathbb{E}_q|e^{j\phi_n}|^2) = \mat{I}$ because the
exponential has unit modulus. The cancellation in
Eq.~\eqref{eq:h-update-3} is what keeps the channel update cheap: the
precision matrix does not depend on the phase posterior at all, so it is
factorised once and reused across iterations.

Holding $q(\vect{h})$ fixed, the phase update is a smoothing problem. The
per-sample log-likelihood contribution is
\begin{equation}
  \label{eq:phase-obs}
  \log p(y_n \mid \phi_n)
  = \frac{2}{\sigma^2} \Re\!\left\{ y_n^{*} e^{j\phi_n} \hat{z}_n \right\}
    + \mathrm{const} ,
  \qquad \hat{z}_n = [\mat{X}\hat{\vect{h}}]_n ,
\end{equation}
which is a von Mises likelihood in $\phi_n$. We approximate it by a Gaussian
matched at its mode, with precision
$\kappa_n = 2 |y_n \hat{z}_n| / \sigma^2$, and then run a Rauch-Tung-Striebel
smoother~\citep{rauch2015smoothing} over the Wiener prior of Eq.~\eqref{eq:phase-prior}. The smoother
returns $\mathbb{E}_q[\phi_n]$ and its variance, from which
$\alpha_n = \exp(-j\mathbb{E}_q[\phi_n] - \tfrac{1}{2}\mathrm{Var}_q[\phi_n])$
follows and closes the loop.

\begin{lemma}[Monotone descent]
\label{lem:descent}
Each pair of updates in Eqs.~\eqref{eq:h-update-4} and~\eqref{eq:phase-obs}
does not increase the variational free energy, and the Gaussian matching step
introduces an error of order $\kappa_n^{-2}$ in the free energy.
\end{lemma}

We initialise $\alpha_n = 1$, which makes the first channel update exactly the
linear minimum mean square error estimate under the common-phase model, so
\est{} begins from the conventional solution and improves it. This also
supplies the graceful degradation observed in Section~\ref{sec:results}: when
the linewidth is negligible the phase posterior stays flat, $\alpha_n$ stays
near unity, and the estimator returns the baseline answer.

\subsection{Termination}

We iterate until the relative change in the free energy falls below $10^{-4}$
or ten iterations elapse. Across the whole sweep in
Section~\ref{sec:results}, the median number of iterations is four and the
ninety-fifth percentile is seven.
